% Abstracts are usually written in English, with a version in your
% mother tongue underneath
\chapter*{Abstract} 
\addcontentsline{toc}{chapter}{Abstract}
% Don't change anything above this.
Malaria remains a major public health challenge in Kenya due to substantial spatial and temporal variation in transmission across different ecological regions. This study developed a spatio-temporal malaria prediction framework to forecast malaria incidence across Kenya’s 47 counties using data collected from January 2016 to December 2025. The framework integrated climatic, environmental, demographic, intervention, and engineered temporal variables, including lag features, rolling averages, and seasonal encodings, to capture delayed climatic effects and seasonal transmission dynamics. Baseline machine learning models, namely Random Forest and XGBoost, were compared with hybrid deep learning architectures comprising CNN-LSTM, CNN-GRU, and CNN-Transformer. Model performance was evaluated using Root Mean Squared Error (RMSE), Mean Absolute Error (MAE), and the coefficient of determination ($R^2$).The results showed clear spatial and temporal variation in malaria transmission across Kenya. XGBoost achieved the highest overall predictive performance (R² = 0.8199), while CNN-LSTM was the best-performing hybrid deep learning model (R² = 0.8096, RMSE = 0.0088, MAE = 0.0052), demonstrating a strong ability to learn complex malaria transmission patterns. SHAP explainability analysis identified the rainfall–NDVI, rainfall–temperature interaction, rainfall, humidity, and vegetation conditions as the most influential predictors of malaria risk. Spatial risk mapping further revealed persistent malaria hotspots in counties surrounding the Lake Victoria basin, including Kisumu, Homa Bay, Migori, and Busia, as well as in coastal counties such as Lamu and Taita Taveta. The findings demonstrate that spatio-temporal machine learning frameworks can provide reliable malaria forecasts and meaningful risk maps for public health planning. The proposed framework offers valuable support for early warning systems, targeted malaria interventions, and evidence-based decision-making aimed at reducing the malaria burden in Kenya.


Keywords: Malaria prediction; Spatio-temporal modelling; Hybrid deep learning; CNN-LSTM; Machine learning; XGBoost; Malaria risk mapping; Kenya.



% At a unviersity like Stellenbosch you *must* produce an abstract in Afrikaans for your masters.
% At AIMS you are encouraged to repeat the abstract in your mother tongue
% French, Igbo, Malagasy, etc. just write it using LaTeX's special
% characters.
% Arabic students use the arabtex package.
% Amharic use openoffice and export from there and import a figure here.
% Where the words do not exist put the English work in italics, or use mathematical symbols.


% Do not change anything below this except for adding your
% signature and name. And take the message to heart.
\vfill
\section*{Declaration}
I, the undersigned, hereby declare that the work contained in this essay is my original work, and that any work done by others or by myself previously has been acknowledged and referenced accordingly.

% Scan your signature into a small picture called 'signature.png' and insert it
% above your name and the date:
\includegraphics[height=4cm]{Figures/My signature.pdf} \hrule
% Your name must be in English Capitalisation with no comma, and the Family name comes last. 
% Do note the date below. It is called the "deadline".
LINDAH CHEROP NGEYWO, 23 June 2026.



